\section{Experiments}
\label{sec:experiments}

We evaluate the framework on three experiments: a toy problem validating gradient correctness, a uniform-coverage baseline that recovers known-optimal geometries, and a weighted target optimization on a problem class where parametric families fail.

\subsection{Common Setup}
\label{sec:common_setup}

All experiments use $\partial$SGP4~\citep{acciarini_closing_2025} for orbit propagation with PyTorch autograd for gradient computation, and AdamW with learning rate $\eta = 10^{-2}$  for optimization.
Circular orbits are assumed and held stationary by zeroing the mean-motion gradient.
The minimum elevation angle for all ground points is $\alpha_{\min} = \SI{10}{\degree}$.
As discussed in \secref{sec:tightness}, we use separate sigmoid temperatures for coverage and revisit: the coverage temperature $\tau_{\text{cov}}$ controls gradient smoothness, while a tighter revisit temperature $\tau_{\text{rev}}$ prevents the leaky integrator from being gamed by phantom coverage.
To select these quantities and the remaining hyperparameters ($\beta$, $\lambda$), we optimize a 24-satellite constellation from four initializations ranging from well-spread to heavily clustered, then grid-search over temperature and weight combinations to find settings that correctly rank the resulting solutions by their hard metrics.
This calibration confirms that $\tau_{\text{cov}} \gg \tau_{\text{rev}}$ is required for correct ranking, though in practice we favor slightly softer values that improve convergence at the expense of absolute correctness (Appendix~\ref{app:hyperparam_grid}).
Coverage and revisit metrics are evaluated over a \SI{24}{\hour} propagation horizon discretized into $K = 240$ time steps on a ground grid of $36 \times 72 = 2592$ points spanning $\pm\SI{70}{\degree}$ latitude.
In each experiment the satellite count $N$ and plane structure are fixed ahead of time and chosen to match the intended comparison geometry (e.g.\ Walker 24/6/1 in \secref{sec:uniform}): this paper addresses the configuration problem, not the sizing problem.
Specific relaxation parameters, optimizer settings, and revisit weight $\lambda$ vary per experiment and are stated in each subsection.

\subsection{Experiment 1: Toy Problem}
\label{sec:toy}
    
To validate that the relaxed objectives produce correct gradients and results, we consider the simplest non-trivial constellation design problem: two satellites in a single orbital plane at \SI{550}{\kilo\metre} altitude with inclination fixed at \SI{60}{\degree} and RAAN fixed at \SI{0}{\degree}, optimizing only their mean anomalies $\mathcal{M}_1$ and $\mathcal{M}_2$.
We use relaxation parameters $\tau_{\text{cov}} = \SI{5}{\degree}$, $\tau_{\text{rev}} = \SI{2}{\degree}$, $\beta = \SI{10}{\minute}$, $\lambda = 2.0$, and AdamW with learning rate $\eta = 10^{-2}$.

The optimal solution is known analytically: the two satellites should be separated by \SI{180}{\degree} in mean anomaly to maximize coverage and revisit rate within the shared orbital plane.
This reduces the problem to two free parameters, allowing the full loss landscape to be visualized.

\paragraph{Loss landscape.}
We evaluate both the relaxed loss and hard (discrete) metrics on a dense grid over $(\mathcal{M}_1, \mathcal{M}_2) \in [0, 360)^2$\SI{}{\degree}.  
\figref{fig:toy_landscape} shows the resulting landscapes.
Both landscapes exhibit a clear global minimum along the anti-diagonal $|\mathcal{M}_1 - \mathcal{M}_2| = \SI{180}{\degree}$, with no local minima.
The relaxed landscape is visibly smoother than the hard landscape while preserving the same basin structure, confirming that the relaxations provide a smooth surface amenable to gradient descent without significantly distorting the underlying objective.

\paragraph{Gradient descent trajectory.}
Starting from $\mathcal{M}_1 = \SI{179}{\degree}$, $\mathcal{M}_2 = \SI{181}{\degree}$ the optimizer converges to $|\mathcal{M}_1 - \mathcal{M}_2| \approx \SI{180}{\degree}$ within 800 iterations.
The trajectory (overlaid in white on \figref{fig:toy_landscape}) follows the gradient field smoothly toward the minimum without oscillation.
This validates that the composed relaxations produce well-behaved gradients---a prerequisite for the higher-dimensional experiments that follow.

\begin{figure}[!htb]
    \centering
    \includegraphics[width=\textwidth]{figures/exp1_gradient_validation.pdf}
    \caption{Two-satellite mean anomaly optimization. \textbf{Top:} relaxed metrics. \textbf{Bottom:} hard metrics. The optimizer (white line) converges to the optimal \SI{180}{\degree} separation. The relaxed landscape preserves the basin structure of the hard landscape while providing smooth gradients.}
    \label{fig:toy_landscape}
\end{figure}

\subsection{Experiment 2: Uniform Coverage Baseline}
\label{sec:uniform}

We now apply our gradient-based scheme to a well-understood uniform coverage problem. We optimize $N = 24$ satellites in 6 orbital planes (4 per plane) with inclination fixed at \SI{60}{\degree}.
RAAN and mean anomaly are free, with RAAN shared within each plane, yielding $6 + 24 = 30$ effective degrees of freedom.

\paragraph{Initial configuration.}
We deliberately initialize the constellation with irregular RAAN offsets $[0, 30, 120, 200, 210, 300]$\SI{}{\degree} far from uniform spacing, and random mean anomalies. The classical Walker-delta $24/6/1$ pattern uses \SI{60}{\degree} RAAN spacing.

\paragraph{Results.}
Over 1000 iterations, the optimizer recovers near-uniform RAAN spacing from the irregular initialization, approaching the Walker-delta geometry without domain-specific inductive biases.
\tabref{tab:uniform_results} summarizes the metrics; \figref{fig:convergence} shows convergence, spatial coverage redistribution, and final RAAN/MA configuration space.

\begin{table}[!htb]
\centering
\caption{Recovery of Walker Delta constellation through gradient optimization of randomly initialized 24-satellite constellation ($6 \times 4$), 1000 iterations, \SI{550}{\kilo\metre} fixed altitude, inclination fixed at \SI{60}{\degree}.}
\label{tab:uniform_results}
\begin{tabular}{lrrr}
\toprule
\textbf{Metric} & \textbf{Initial} & \textbf{Optimized} & \textbf{Walker 24/6/1} \\
\midrule
Hard coverage [\%] & 32.41 & 40.34 & 40.33 \\
Hard mean max revisit [min] & 136.2 & 48.0 & 48.0 \\
\midrule
Soft coverage [\%] & 41.86 & 50.32 & 50.32 \\
Soft mean max revisit [min] & 127.5 & 73.5 & 73.5 \\
\bottomrule
\end{tabular}
\end{table}

\begin{figure}[!htb]
    \centering
    \begin{minipage}[t]{0.48\textwidth}
        \centering
        \includegraphics[width=\textwidth]{figures/exp2_convergence.pdf}\\[-2pt]
        {\small (a) Convergence}
    \end{minipage}%
    \hfill
    \begin{minipage}[t]{0.48\textwidth}
        \centering
        \includegraphics[width=\textwidth]{figures/exp2_raan_ma.pdf}\\[-2pt]
        {\small (b) RAAN vs mean anomaly}
    \end{minipage}
    \caption{Walker recovery experiment. \textbf{(a)} Convergence of coverage and revisit over 1000 iterations, with Walker 24/6/1 reference (dashed). \textbf{(b)} RAAN vs mean anomaly trajectories (hollow = initial, solid = optimized, stars = Walker). The optimizer recovers Walker-equivalent geometry (in rotational symmetry) from an irregular initialization.}
    \label{fig:convergence}
\end{figure}

\paragraph{Loss landscape.}
To investigate the structure of the loss landscape, we visualize it using per-trajectory PCA~\citep{li_visualizing_2018}: for each of the four initializations from the hyperparameter calibration (Appendix~\ref{app:hyperparam_grid}), we project the optimization trajectory onto its two principal components and evaluate the loss on that 2D slice.
\figref{fig:landscape} shows smooth basins in all four cases, confirming that the relaxations produce a landscape amenable to gradient descent. However, the global view reveals multi-modal structure: the clustered initializations converge to distinct local minima with worse hard metrics than the well-spread ones, indicating that initialization still matters despite the smooth local landscape.

\begin{figure*}[!htb]
    \centering
    \includegraphics[width=\textwidth]{figures/exp3_loss_landscape.pdf}
    \caption{Loss landscape visualization~\citep{li_visualizing_2018}. \textbf{Left:} global view along random directions. \textbf{Columns 2--5:} per-trajectory PCA slices (white line = optimizer path). \textbf{Top/bottom:} relaxed and hard loss, respectively. Per-trajectory views show smooth basins; the global view reveals multi-modal structure.}
    \label{fig:landscape}
\end{figure*}

\subsection{Experiment 3: Regional Target Optimization}
\label{sec:weighted}

The primary advantage of gradient-based constellation optimization over parametric families is the ability to handle non-uniform, mission-specific objectives.
We demonstrate this by optimizing a small constellation for coverage and revisit over Europe---a problem that Walker constellations cannot address by construction.

\paragraph{Setup.}
We sample 500 ground targets from European country boundaries (GeoJSON), weighted by $\cos(\text{lat})$ to correct for area distortion.
The constellation consists of 4 satellites in 2 orbital planes (2 per plane).
Unlike Experiment~2, we free all orbital geometry parameters: inclination ($\iota \in [30, 90]$\SI{}{\degree}), RAAN, mean anomaly, argument of perigee, and orbital shape.
To allow the optimizer to explore eccentric orbits without producing unphysical perigee altitudes, we reparameterize eccentricity and altitude jointly as (perigee altitude, excess altitude), where perigee $\in [400, 600]$\SI{}{\kilo\metre} and excess $\in [0, 1500]$\SI{}{\kilo\metre}, guaranteeing perigee $\geq$ \SI{400}{\kilo\metre} by construction.
Inclination, RAAN, eccentricity, argument of perigee, and altitude are shared within each plane; mean anomaly is free per satellite, giving 12 effective degrees of freedom.
We use $\lambda = 1.0$, and run for 3000 iterations.

\paragraph{Results.}
Starting from a Walker-like circular initialization at \SI{550}{\kilo\metre}, the optimizer discovers a qualitatively different geometry (\tabref{tab:weighted_results}).
Both planes converge to highly elliptical orbits ($e \approx 0.40$) with argument of perigee $\omega \approx \SI{270}{\degree}$, placing apogee over the northern hemisphere where the satellite dwells longest.
The resulting orbits have perigee $\approx$ \SI{685}{\kilo\metre} and apogee $\approx$ \SI{10000}{\kilo\metre}---a Molniya-class highly elliptical orbit discovered entirely from gradients, with no domain-specific heuristics.
The inclination converges to $\sim$\SI{63}{\degree}--\SI{65}{\degree}, near the critical inclination (\SI{63.4}{\degree}) at which $J_2$ perturbations do not precess the argument of perigee.

The optimized constellation achieves 99.3\% coverage and 3.7~min mean revisit over Europe, compared to 10.0\% coverage and 247.5~min for a Walker 4/2/1 baseline (\tabref{tab:weighted_results}).
The comparison is not between equal orbit regimes: the Walker baseline is constrained to circular LEO while the optimizer is free to explore eccentric orbits. The result demonstrates that, given this freedom, the optimizer escapes the LEO basin and discovers a qualitatively different and far more effective geometry for regional coverage.

\figref{fig:europe_density} shows the visibility density (number of timesteps each ground bin is within \SI{10}{\degree} elevation of any satellite).
The optimized constellation concentrates nearly all coverage in the northern hemisphere over Europe, while the Walker baseline distributes coverage uniformly but achieves far less density.
The orbital period of $\sim$\SI{210}{\minute} produces approximately 7 revolutions per sidereal day, creating a near-resonant repeating ground track.

\begin{table}[!htb]
\centering
\caption{Regional target optimization: 4 satellites (2 planes $\times$ 2), 3000 iterations, all orbital geometry parameters free with coupled perigee/apoapsis constraints.}
\label{tab:weighted_results}
\begin{tabular}{lrr}
\toprule
\textbf{Metric} & \textbf{Walker 4/2/1} & \textbf{Optimized} \\
\midrule
Hard coverage [\%] & 9.97 & 99.25 \\
Hard mean max revisit [min] & 247.5 & 3.7 \\
Soft coverage [\%] & 12.37 & 97.99 \\
Soft mean max revisit [min] & 208.9 & 54.9 \\
\bottomrule
\end{tabular}
\end{table}

\begin{table}[!htb]
\centering
\caption{Per-plane orbital geometry discovered by the optimizer for the Europe target set. Both planes converge to eccentric, Molniya-like orbits with apogee over the northern hemisphere ($\omega \approx \SI{270}{\degree}$). The Walker baseline uses circular orbits at \SI{550}{\kilo\metre}.}
\label{tab:exp3_geometry}
\begin{tabular}{lrrrrr}
\toprule
& $\iota$ \textbf{[deg]} & $\Omega$ \textbf{[deg]} & $e$ & $\omega$ \textbf{[deg]} & \textbf{Perigee / Apogee [km]} \\
\midrule
Plane 0 & 62.6 & 336.6 & 0.398 & 276.2 & 685 / 10011 \\
Plane 1 & 65.3 & 155.6 & 0.402 & 266.5 & 687 / 10193 \\
\midrule
Walker (all) & 60.0 & uniform & 0.001 & --- & 550 / 550 \\
\bottomrule
\end{tabular}
\end{table}

\begin{figure}[htbp]
    \centering
    \includegraphics[width=0.48\textwidth]{figures/exp3_convergence.pdf}
    \caption{Convergence of the Europe target optimization over 3000 iterations. The optimizer rapidly improves both coverage and revisit over the Walker baseline (dashed).}
    \label{fig:exp3_convergence}
\end{figure}

\begin{figure}[!htb]
    \centering
    \includegraphics[width=.9\textwidth]{figures/exp3_groundtrack_optimized.pdf}\\[2pt]
    {\small (a) Optimized for Europe}\\[8pt]
    \includegraphics[width=.9\textwidth]{figures/exp3_groundtrack_walker.pdf}\\[2pt]
    {\small (b) Walker 4/2/1 baseline}
    \caption{Visibility density over \SI{24}{\hour}. \textbf{(a)} The optimized constellation concentrates coverage over Europe via eccentric orbits with apogee dwell over northern latitudes. \textbf{(b)} The Walker baseline distributes coverage uniformly but achieves far less density over the target region.}
    \label{fig:europe_density}
\end{figure}

\subsection{Comparison to Baseline Methods}
\label{sec:baselines}

Constellation design in prior work is typically posed as a black-box optimisation over design variables, handled either by population heuristics (simulated annealing, genetic algorithms) or off-the-shelf gradient-free libraries.
We benchmark our differentiable pipeline against three such baselines on the Walker recovery testbed of \secref{sec:uniform}: 24 satellites in a $6 \times 4$ arrangement at \SI{550}{\kilo\metre}, inclination fixed at \SI{60}{\degree}, and the same irregular-RAAN initialisation.
The decision vector has 30 continuous variables (six per-plane RAANs and 24 per-satellite mean anomalies).

\paragraph{Baselines.}
We compare against simulated annealing (SA), a genetic algorithm (GA), and differential evolution (DE), all applied to the same fitness $J = -C + \lambda \Delta^{\max}$ (with $\lambda = 0.1$ as in \secref{sec:uniform}), computed directly from the \emph{hard} coverage and revisit metrics of \eqref{eq:coverage_fraction_hard} and \eqref{eq:mean_revisit_hard}---the baselines do not see the relaxed objective.
SA and GA follow the constellation-design formulations used in prior heuristic work~\citep{paek_optimization_2019} with standard tuning: SA uses a probe-calibrated temperature schedule (\SI{80}{\percent} initial / \SI{1}{\percent} final acceptance, geometric cooling) and an adaptive step size targeting the Vanderbilt--Louie rate of $0.44$; GA uses tournament selection (size 3), uniform crossover, and Gaussian mutation with $\sigma$ annealed from $\SI{30}{\degree}$ to $\SI{5}{\degree}$.
DE uses the TwoPointsDE implementation from nevergrad~\citep{rapin_nevergrad_2018}.
All three baselines warm-start from the same irregular-RAAN configuration as the gradient run.
Each is run for 5 seeds with a budget of 4050 function evaluations (4000 main evaluations plus the 50-evaluation probe that SA uses to calibrate its temperature schedule, matched across methods); the gradient pipeline uses 1000 iterations of the relaxed objective.

\begin{table}[!htb]
\centering
\caption{Final hard metrics on the Walker recovery testbed. The gradient pipeline recovers the Walker reference exactly in 1000 iterations; the three black-box baselines plateau well short of it even with roughly four times the evaluation budget and best-of-5 seed selection.}
\label{tab:baselines}
% Manually generated from experiments/exp2_baselines.py (5 seeds, 4050 evals each: 4000 main + 50-eval SA probe, matched across methods)
% Range is min--max across seeds; the best side of each range is bolded
% (highest coverage, lowest revisit).
\begin{tabular}{lrr}
\toprule
\textbf{Method} & \textbf{Hard coverage [\%]} & \textbf{Hard mean worst-case revisit [min]} \\
\midrule
Walker 24/6/1 reference & 40.33 & 48.0 \\
Gradient (ours, 1000 iters) & \textbf{40.34} & \textbf{48.0} \\
\midrule
\multicolumn{3}{l}{\emph{Black-box baselines (4000 iters, range over 5 seeds):}} \\
\quad Simulated annealing       & 34.33--\textbf{36.54} & \textbf{73.9}--80.5 \\
\quad Genetic algorithm          & 37.83--\textbf{38.69} & \textbf{60.3}--62.2 \\
\quad TwoPointsDE~\citep{rapin_nevergrad_2018} & 35.77--\textbf{37.44} & \textbf{68.5}--76.1 \\
\bottomrule
\end{tabular}

\end{table}

\figref{fig:baselines_conv} plots convergence for the best-fitness seed of each method against function evaluations; the $x = 0$ values carry the common bad initialisation forward by zero-order hold.
The differentiable pipeline reaches the Walker reference (\SI{40.3}{\percent} coverage, \SI{48}{\minute} revisit) at roughly 750 evaluations and is flat thereafter.
Among the heuristics, GA performs best (\SI{60.3}{\min} revisit, best seed), followed by DE (\SI{68.5}{\min}) and SA (\SI{73.9}{\min}); in our experiments with these tuned baselines, none come within \SI{12}{\minute} of Walker revisit within the \num{4050}-evaluation budget.
GA's edge over DE on this problem is attributable to its uniform crossover and $\bmod\,\SI{360}{\degree}$ wrap, which respect the cyclic nature of RAAN and mean anomaly; DE's Euclidean-space perturbations do not.\footnote{The nevergrad NGOpt meta-selector chooses CMA-ES with box constraints at this budget and dimension; we found it plateaus earlier than TwoPointsDE and report the stronger of the two.}
The gap between the gradient pipeline and all three baselines is consistent with the information-per-evaluation argument: each gradient step uses a 30-dimensional descent direction obtained by autodiff through the relaxed objective, whereas the baselines receive a single scalar per evaluation.

\begin{figure}[!htb]
    \centering
    \begin{minipage}[t]{0.49\linewidth}\centering
        \includegraphics[width=\textwidth]{figures/exp2_method_comparison_coverage.pdf}\\[-2pt]
        {\small (a) Hard coverage}
    \end{minipage}\hfill
    \begin{minipage}[t]{0.49\linewidth}\centering
        \includegraphics[width=\textwidth]{figures/exp2_method_comparison_revisit.pdf}\\[-2pt]
        {\small (b) Hard mean worst-case revisit}
    \end{minipage}
    \caption{Convergence of the differentiable optimiser (teal, Exp 2 baseline) against three black-box baselines on the Walker recovery testbed. Best-of-5 seeds shown per method; evaluation $0$ shares the common irregular-RAAN initialisation by zero-order hold, and the Walker reference is dotted grey. The gradient pipeline matches Walker at ${\sim}750$ evaluations; the heuristics plateau above Walker revisit even with roughly four times the budget.}
    \label{fig:baselines_conv}
\end{figure}

\subsection{Ablations}
\label{sec:ablations}

To isolate the design choices behind the relaxed pipeline, we re-run the Walker recovery testbed of \secref{sec:uniform} (24 satellites, $6 \times 4$, inclination fixed at \SI{60}{\degree}, RAAN and mean anomaly free, irregular RAAN init, 1000 iterations) under five controlled perturbations.
Constellation, ground grid, propagation horizon, and optimizer are held fixed; only the ablated knob varies.
Each row of \tabref{tab:ablations} reports the final hard coverage and mean worst-case revisit; the Walker 24/6/1 reference is included for calibration.

\begin{table}[!htb]
\centering
\caption{Design-choice ablations on the Walker recovery testbed (\secref{sec:uniform}, 24 satellites, 1000 iterations, hard metrics evaluated at the final step). Default settings (used elsewhere unless ablated) are $\lambda = 0.1$, $\tau_{\text{cov}} = \tau_{\text{rev}} = \SI{2}{\degree}$, $\beta = \SI{10}{\minute}$.}
\label{tab:ablations}
% Auto-generated by experiments/exp4_ablations.py
\begin{tabular}{lrr}
\toprule
\textbf{Setting} & \textbf{Hard coverage [\%]} & \textbf{Hard mean worst-case revisit [min]} \\
\midrule
Walker 24/6/1 reference / Exp 2 defaults & 40.34 & 48.0 \\
\midrule
(a) Coverage only ($\lambda = 0$) & 40.76 & 126.7 \\
(a) Revisit only ($-\tilde{C}$ dropped) & 40.34 & 48.0 \\
(b) Split $\tau$ ($\tau_{\text{cov}}=3^{\circ}$, $\tau_{\text{rev}}=1^{\circ}$) & 39.04 & 54.9 \\
(b) Split $\tau$ inverted ($\tau_{\text{cov}}=1^{\circ}$, $\tau_{\text{rev}}=3^{\circ}$) & 40.11 & 58.3 \\
(c) $\lambda = 0.01$ & 40.06 & 59.8 \\
(c) $\lambda = 1.0$ & 40.34 & 47.9 \\
(d) $\beta = 1$ min & 38.58 & 55.3 \\
(d) $\beta = 100$ min & 40.08 & 68.9 \\
(e) Random RAAN+MA init, $n=10$ & 39.73 $\pm$ 0.46 (38.82--40.34) & 61.7 $\pm$ 4.7 (47.8--64.1) \\
\bottomrule
\end{tabular}

\end{table}

\paragraph{(a) Loss composition.}
Coverage-only preserves Walker-level coverage (\SI{40.8}{\percent}) but leaves revisit at \SI{126.7}{\minute}---over \SI{2.6}{\times} the Walker baseline---because the loss never penalises long gaps.
In contrast, revisit-only and the full combined loss both converge to the Walker-equivalent solution (\SI{40.3}{\percent} coverage, \SI{48.0}{\minute} revisit): once satellites are spaced to minimise revisit, the coverage term contributes only a small residual gradient, so the two are effectively equivalent on this testbed.
The combined form remains load-bearing on problems where clustering is a valid local minimum of the revisit objective (\secref{sec:weighted}); here it confirms that adding the coverage term does no harm.

\paragraph{(b) Sigmoid temperature.}
We compare the shared baseline ($\tau_{\text{cov}} = \tau_{\text{rev}} = \SI{2}{\degree}$, used in \secref{sec:uniform}) against the split values calibrated in Appendix~\ref{app:hyperparam_grid} ($\tau_{\text{cov}} = \SI{3}{\degree}$, $\tau_{\text{rev}} = \SI{1}{\degree}$) and the inverse split ($\tau_{\text{cov}} = \SI{1}{\degree}$, $\tau_{\text{rev}} = \SI{3}{\degree}$).
The shared setting recovers Walker (\SI{40.3}{\percent}, \SI{48.0}{\minute}); both split variants underperform (\SI{39.0}{\percent} / \SI{54.9}{\minute} and \SI{40.1}{\percent} / \SI{58.3}{\minute}, respectively).
The split temperatures' role, as discussed in \secref{sec:relaxation_quality}, is to preserve correct \emph{ranking} of distinct converged solutions in the multi-modal regime; for optimising a single trajectory from a reasonable initialisation, shared temperature is both simpler and better.

\paragraph{(c) Revisit weight $\lambda$.}
Sweeping $\lambda \in \{0.01, 0.1, 1.0\}$ shows that results are insensitive to $\lambda$ over two orders of magnitude: both $\lambda = 0.1$ and $\lambda = 1.0$ recover Walker (\SI{40.3}{\percent}, \SI{48.0}{\minute}).
Under-weighting revisit at $\lambda = 0.01$ lets coverage saturate before the revisit gradient does useful work (\SI{40.1}{\percent}, \SI{59.8}{\minute}).
$\lambda = 0.1$ is a natural default; the test suggests a broad plateau around it rather than a sharp optimum.

\paragraph{(d) LogSumExp temperature $\beta$.}
Sweeping $\beta \in \{1, 10, 100\}$ minutes exposes both ends of the sharpness/smoothness trade-off.
The tight $\beta = \SI{1}{\minute}$ approaches a hard max and starves near-worst timesteps of gradient, degrading coverage to \SI{38.6}{\percent}.
The loose $\beta = \SI{100}{\minute}$ over-smooths the worst-case signal so the optimiser underweights long gaps, raising revisit to \SI{68.9}{\minute}.
$\beta = \SI{10}{\minute}$ sits between these failure modes and recovers Walker exactly.

\paragraph{(e) Initialization sensitivity.}
Re-running the same configuration under $n = 10$ seeds with \emph{both} RAAN and mean anomaly drawn uniformly at random yields hard coverage in $[38.8, 40.3]\,\%$ and revisit in $[47.8, 64.1]$~min, with mean $39.7 \pm 0.5\,\%$ and $61.7 \pm 4.7$~min.
Only one of the ten seeds recovers Walker exactly; the other nine converge to a worse local minimum near \SI{63}{\minute} revisit while still reaching roughly Walker-level coverage.
This bimodality is consistent with the multi-modal structure visible in \figref{fig:landscape}: coverage drives satellites apart to near-uniform angular spacing reliably, but the finer-grained plane-phasing that distinguishes the global Walker minimum from nearby local minima is not always found from arbitrary RAAN+MA initialisations.
The result of \secref{sec:uniform} is therefore not guaranteed from every initialisation, but remains robustly reachable when RAAN is seeded near a Walker-like spacing (as in all preceding experiments).
A simple best-of-$n$ multi-start ensemble recovers Walker exactly here (seed 5: \SI{40.3}{\percent}, \SI{47.8}{\minute}), so the multi-modality can be mitigated at low cost given parallel runs.

\begin{figure}[!htb]
    \centering
    \begin{minipage}[t]{0.32\textwidth}\centering
        \includegraphics[width=\textwidth]{figures/ablations_loss_revisit.pdf}\\[-2pt]
        {\small (a) Loss composition: revisit}
    \end{minipage}\hfill
    \begin{minipage}[t]{0.32\textwidth}\centering
        \includegraphics[width=\textwidth]{figures/ablations_lambda_revisit.pdf}\\[-2pt]
        {\small (b) Revisit weight $\lambda$: revisit}
    \end{minipage}\hfill
    \begin{minipage}[t]{0.32\textwidth}\centering
        \includegraphics[width=\textwidth]{figures/ablations_beta_revisit.pdf}\\[-2pt]
        {\small (c) LSE temperature $\beta$: revisit}
    \end{minipage}
    \\[6pt]
    \begin{minipage}[t]{0.32\textwidth}\centering
        \includegraphics[width=\textwidth]{figures/ablations_loss_coverage.pdf}\\[-2pt]
        {\small (d) Loss composition: coverage}
    \end{minipage}\hfill
    \begin{minipage}[t]{0.32\textwidth}\centering
        \includegraphics[width=\textwidth]{figures/ablations_lambda_coverage.pdf}\\[-2pt]
        {\small (e) Revisit weight $\lambda$: coverage}
    \end{minipage}\hfill
    \begin{minipage}[t]{0.32\textwidth}\centering
        \includegraphics[width=\textwidth]{figures/ablations_beta_coverage.pdf}\\[-2pt]
        {\small (f) LSE temperature $\beta$: coverage}
    \end{minipage}
    \caption{Hard metrics over training for the three multi-variant ablations; Walker 24/6/1 reference dashed. Top row: mean worst-case revisit. Bottom row: coverage. \textbf{(a, d)} Coverage-only flattens coverage near Walker immediately but never reduces revisit; revisit-only and combined both converge to the Walker solution. \textbf{(b, e)} $\lambda = 0.1$ and $\lambda = 1.0$ are indistinguishable on this testbed; $\lambda = 0.01$ plateaus early on revisit. \textbf{(c, f)} $\beta = \SI{10}{\minute}$ converges cleanly; $\beta = \SI{1}{\minute}$ is gradient-starved on revisit and leaves coverage short; $\beta = \SI{100}{\minute}$ over-smooths and gets stuck at a worse revisit minimum.}
    \label{fig:ablations_conv}
\end{figure}

% \subsection{Computational Cost}
% \label{sec:cost}

% The per-iteration cost scales as $\mathcal{O}(NKP)$---linearly in constellation size $N$, number of time steps $K$, and number of targets $P$.
% For $N = 12$, $K = 200$, $P = 648$, a single iteration takes approximately \SI{2}{\second} on a CPU (Apple M3), and the full 300-iteration optimization completes in $\sim$\SI{10}{\minute}.

% For comparison, \citet{paek_optimization_2019} required $\sim$250 SA iterations or $\sim$2500 GA evaluations, each involving a full STK orbit simulation substantially more expensive than our PyTorch forward pass.
% The gradient-based approach thus achieves comparable or better solution quality with orders-of-magnitude fewer expensive evaluations, and each evaluation is itself cheaper due to the lightweight dSGP4 propagator.
% GPU acceleration via PyTorch batched operations could reduce per-iteration cost further for larger constellations.
